\documentclass[12pt]{article}
\usepackage[a3paper, margin=1in]{geometry}
\usepackage{amsmath, amssymb, amsthm, ulem}

\begin{document}

\section*{Domácí úkol 4.2}

Vypracoval: Daniel \textit{"Randál"} Ransdorf \hfill Podpis: \rule{4cm}{0.4pt}

\begin{enumerate}
  \setcounter{enumi}{1}
  \item \textbf{Řešení}:
    Kolmou projekci vektoru $\vec{x}$ na přímku v trojsložkovém prostoru sestrojíme
    přičtením lineární kombinace dvou různoběžných vektorů, které jsou kolmé na danou přímku.
    Vektory musí tvořit rovinu kolmou na přímku a na té rovině posuneme $\vec{x}$ na přímku.
    Zvolme tedy dva zjevně nezávislé vektory, jejichž skalární součin se směrovým vektorem přímky vychází 0 (kolmost na přímku)
    ... $(-2, 1, 0)^T, (0, 1, -2)^T$

    Pro zobrazení vektoru $\vec{x}$ stačí vyřešit tuto rovnici v realných číslech:
    
    \begin{align*}
      \vec{x} + a \left(\begin{array}{c} -2 \\ 1 \\ 0 \end{array}\right)
      + b \left(\begin{array}{c} 0 \\ 1 \\ -2 \end{array}\right)
      &= t \left(\begin{array}{c} 1 \\ 2 \\ 1 \end{array}\right) \\
      \vec{x} &= t \left(\begin{array}{c} 1 \\ 2 \\ 1 \end{array}\right)
      - a \left(\begin{array}{c} -2 \\ 1 \\ 0 \end{array}\right)
      - b \left(\begin{array}{c} 0 \\ 1 \\ -2 \end{array}\right) \\
      \vec{x} &= \left(\begin{array}{ccc} 1 & -2 & 0 \\ 2 & 1 & 1 \\ 1 & 0 & -2 \end{array}\right)
      \cdot \left(\begin{array}{c} t \\ -a \\ -b \end{array}\right)
    \end{align*}

    Pro nalezení matice A stačí najít obraz trojsložkových kanonických bázových vektorů.
    Obrazy těchto vektorů jsou sloupcemi matice A (důkaz níže).
    Nalezněme tedy pro kanonické bázové vektory parametr $t$, dle kterého poté najdeme obraz vektorů.

    \begin{align*}
      &\left(\begin{array}{ccc} 1 & -2 & 0 \\ 2 & 1 & 1 \\ 1 & 0 & -2 \end{array}\right)
      \cdot \left(\begin{array}{c} t \\ -a \\ -b \end{array}\right) = \left(\begin{array}{c} 1 \\ 0 \\ 0 \end{array}\right)
      \\
      &\left(\begin{array}{ccc} 1 & -2 & 0 \\ 2 & 1 & 1 \\ 1 & 0 & -2 \end{array}\right)
      \cdot \left(\begin{array}{c} t \\ -a \\ -b \end{array}\right) = \left(\begin{array}{c} 0 \\ 1 \\ 0 \end{array}\right)
      \\
      &\underbrace{\left(\begin{array}{ccc} 1 & -2 & 0 \\ 2 & 1 & 1 \\ 1 & 0 & -2 \end{array}\right)
      \cdot \left(\begin{array}{c} t \\ -a \\ -b \end{array}\right) = \left(\begin{array}{c} 0 \\ 0 \\ 1 \end{array}\right)
      } \\
      \text{Použijme soustavy najednou pomocí}&\text{ násobení matic a ignorujme proměnné $a,b$, které nepotřebujeme vyřešit} \\
      &\left(\begin{array}{ccc} 1 & -2 & 0 \\ 2 & 1 & 1 \\ 1 & 0 & -2 \end{array}\right)
      \cdot \left(\begin{array}{ccc} t_1 & t_2 & t_3 \\ ... & ... & ... \\ ... & ... & ... \end{array}\right) = \left(\begin{array}{ccc} 1 & 0 & 0 \\ 0 & 1 & 0 \\ 0 & 0 & 1 \end{array}\right)
    \end{align*}

    \[
      \left(\begin{array}{ccc|ccc} 
        1 & -2 & 0 & 1 & 0 & 0 \\
        2 & 1 & 1 & 0 & 1 & 0\\
        1 & 0 & -2 & 0 & 0 & 1 
      \end{array}\right)
      \overset{\text{3.ř } - \text{1.ř}}{\;\sim\;}
      \left(\begin{array}{ccc|ccc} 
        1 & -2 & 0 & 1 & 0 & 0 \\
        2 & 1 & 1 & 0 & 1 & 0\\
        0 & 2 & -2 & -1 & 0 & 1 
      \end{array}\right)
      \overset{\text{1.ř } \cdot(-2)}{\;\sim\;}
      \left(\begin{array}{ccc|ccc} 
        -2 & 4 & 0 & -2 & 0 & 0 \\
        2 & 1 & 1 & 0 & 1 & 0\\
        0 & 2 & -2 & -1 & 0 & 1 
      \end{array}\right)
      \overset{\text{2.ř } +\text{1.ř}}{\;\sim\;}
      \left(\begin{array}{ccc|ccc} 
        -2 & 4 & 0 & -2 & 0 & 0 \\
        0 & 5 & 1 & -2 & 1 & 0\\
        0 & 2 & -2 & -1 & 0 & 1 
      \end{array}\right)
    \]
    \renewcommand{\arraystretch}{1.2}
    \[
      \overset{\text{2.ř } \cdot -\frac{2}{5}}{\;\sim\;}
      \left(\begin{array}{ccc|ccc} 
        -2 & 4 & 0 & -2 & 0 & 0 \\
        0 & -2 & -\frac{2}{5} & \frac{4}{5} & -\frac{2}{5} & 0 \\
        0 & 2 & -2 & -1 & 0 & 1 
      \end{array}\right)
      \overset{\text{3.ř } +\text{2.ř}}{\;\sim\;}
      \left(\begin{array}{ccc|ccc} 
        -2 & 4 & 0 & -2 & 0 & 0 \\
        0 & -2 & -\frac{2}{5} & \frac{4}{5} & -\frac{2}{5} & 0 \\
        0 & 0 & -\frac{12}{5} & -\frac{1}{5} & -\frac{2}{5} & 1
      \end{array}\right)
      \overset{\text{3.ř } \cdot -\frac{1}{6}}{\;\sim\;}
      \left(\begin{array}{ccc|ccc} 
        -2 & 4 & 0 & -2 & 0 & 0 \\
        0 & -2 & -\frac{2}{5} & \frac{4}{5} & -\frac{2}{5} & 0 \\
        0 & 0 & \frac{2}{5} & \frac{1}{30} & \frac{1}{15} & -\frac{1}{6}
      \end{array}\right)
    \]
    \[
      \overset{\text{2.ř } +\text{3.ř}}{\;\sim\;}
      \left(\begin{array}{ccc|ccc} 
        -2 & 4 & 0 & -2 & 0 & 0 \\
        0 & -2 & 0 & \frac{5}{6} & -\frac{1}{3} & -\frac{1}{6} \\
        0 & 0 & \frac{2}{5} & \frac{1}{30} & \frac{1}{15} & -\frac{1}{6}
      \end{array}\right)
      \overset{\text{2.ř } \cdot 2}{\;\sim\;}
      \left(\begin{array}{ccc|ccc} 
        -2 & 4 & 0 & -2 & 0 & 0 \\
        0 & -4 & 0 & \frac{5}{3} & -\frac{2}{3} & -\frac{1}{3} \\
        0 & 0 & \frac{2}{5} & \frac{1}{30} & \frac{1}{15} & -\frac{1}{6}
      \end{array}\right)
      \overset{\text{1.ř } +\text{2.ř}}{\;\sim\;}
      \left(\begin{array}{ccc|ccc} 
        -2 & 0 & 0 & -\frac{1}{3} & -\frac{2}{3} & -\frac{1}{3} \\
        0 & -4 & 0 & \frac{5}{3} & -\frac{2}{3} & -\frac{1}{3} \\
        0 & 0 & \frac{2}{5} & \frac{1}{30} & \frac{1}{15} & -\frac{1}{6}
      \end{array}\right)
    \]
    \[
      \overset{\text{1.ř } \cdot -\frac{1}{2}}{\;\sim\;}
      \left(\begin{array}{ccc|ccc} 
        1 & 0 & 0 & \frac{1}{6} & \frac{1}{3} & \frac{1}{6} \\
        0 & -4 & 0 & \frac{5}{3} & -\frac{2}{3} & -\frac{1}{3} \\
        0 & 0 & \frac{2}{5} & \frac{1}{30} & \frac{1}{15} & -\frac{1}{6}
      \end{array}\right)
    \]

    Na prvním řádku na pravé straně máme chtěné hodnoty $t_1, t_2, t_3$.
    Položme $A = \left(\begin{array}{c} \vec{a_1} \;\vline\; \vec{a_2} \;\vline\; \vec{a_3} \end{array}\right)$,
    a ukažme, že sloupcové vektory jsou rovny obrazům kanonických bázových vektorů, které získáme dosazením parametrů $t_1, t_2, t_3$.

    \renewcommand{\arraystretch}{1.3}
    \begin{align*}
      \vec{a_1} = \vec{a_1} + 0 + 0 = A \cdot \left(\begin{array}{c} 1 \\ 0 \\ 0 \end{array}\right)
        = t_1 \left(\begin{array}{c} 1 \\ 2 \\ 1 \end{array}\right)
        = \frac{1}{6} \left(\begin{array}{c} 1 \\ 2 \\ 1 \end{array}\right)
        = \left(\begin{array}{c} \frac{1}{6} \\ \frac{1}{3} \\ \frac{1}{6} \end{array}\right) \\
      \vec{a_2} = 0 + \vec{a_2} + 0 = A \cdot \left(\begin{array}{c} 0 \\ 1 \\ 0 \end{array}\right)
        = t_2 \left(\begin{array}{c} 1 \\ 2 \\ 1 \end{array}\right)
        = \frac{1}{3} \left(\begin{array}{c} 1 \\ 2 \\ 1 \end{array}\right)
        = \left(\begin{array}{c} \frac{1}{3} \\ \frac{2}{3} \\ \frac{1}{3} \end{array}\right) \\
      \vec{a_3} = 0 + 0 + \vec{a_3} = A \cdot \left(\begin{array}{c} 0 \\ 0 \\ 1 \end{array}\right)
        = t_3 \left(\begin{array}{c} 1 \\ 2 \\ 1 \end{array}\right)
        = \frac{1}{6} \left(\begin{array}{c} 1 \\ 2 \\ 1 \end{array}\right)
        = \left(\begin{array}{c} \frac{1}{6} \\ \frac{1}{3} \\ \frac{1}{6} \end{array}\right) \\
    \end{align*}

    Dosazme sloupcové vektory a získejme matici A:

    \[
      A = \left(\begin{array}{c} \vec{a_1} \;\vline\; \vec{a_2} \;\vline\; \vec{a_3} \end{array}\right)
        = \underline{\underline{\left(\begin{array}{ccc}
            \frac{1}{6} & \frac{1}{3} & \frac{1}{6} \\
            \frac{1}{3} & \frac{2}{3} & \frac{1}{3} \\
            \frac{1}{6} & \frac{1}{3} & \frac{1}{6} \\
          \end{array}\right)}}
    \]

\end{enumerate}

\textbf{BONUSOVÝ PROBLÉM:}
  
Úlohu řešme podobně. Zparametrizujme si nejprve zadanou rovinu, položme tedy $y=t_1, z=t_2$:
    \begin{align*}
      ax + bt_1 + ct_2 &= 0 \\
      ax &= -bt_1 - ct_2 \\
      x &= -\frac{b}{a}t_1 - \frac{c}{a}t_2 \\
      \Longrightarrow \left(\begin{array}{c} x\\y\\z \end{array}\right)
        &= t_1 \left(\begin{array}{c} -\frac{b}{a}\\1\\0 \end{array}\right) + t_2 \left(\begin{array}{c} -\frac{c}{a}\\0\\1 \end{array}\right)
    \end{align*}

  Abychom zadaným směrovým vektorem posunuli vektor $\vec{v}$ na rovinu, musíme vyřešit rovnici:

    \begin{align*}
      \vec{v} + k \left(\begin{array}{c} d\\e\\f \end{array}\right) &=
        t_1 \left(\begin{array}{c} -\frac{b}{a}\\1\\0 \end{array}\right) + t_2 \left(\begin{array}{c} -\frac{c}{a}\\0\\1 \end{array}\right) \\
      \vec{v} &=
        t_1 \left(\begin{array}{c} -\frac{b}{a}\\1\\0 \end{array}\right) + t_2 \left(\begin{array}{c} -\frac{c}{a}\\0\\1 \end{array}\right)
        - k \left(\begin{array}{c} d\\e\\f \end{array}\right) \\
      \vec{v} &=
        \left(\begin{array}{ccc} -\frac{b}{a} & -\frac{c}{a} & d \\ 1 & 0 & e \\ 0 & 1 & f \end{array}\right) \left(\begin{array}{c} t_1\\t_2\\-k \end{array}\right)
    \end{align*}
  
  Opět vyřešme pro kanonické bázové vektory:

    \[
      \left(\begin{array}{ccc|ccc} 
        -\frac{b}{a} & -\frac{c}{a} & d & 1 & 0 & 0 \\
        1 & 0 & e & 0 & 1 & 0\\
        0 & 1 & f & 0 & 0 & 1 
      \end{array}\right)
      \overset{\text{Eliminace Wolframem}}{\;\sim\;}
      \left(
      \begin{array}{ccc|ccc}
      1 & 0 & 0 &
      -\frac{a e}{a d + b e + c f} &
      \frac{a d + c f}{a d + b e + c f} &
      -\frac{c e}{a d + b e + c f} \\
      0 & 1 & 0 &
      -\frac{a f}{a d + b e + c f} &
      -\frac{b f}{a d + b e + c f} &
      \frac{a d + b e}{a d + b e + c f} \\
      0 & 0 & 1 &
      \frac{a}{a d + b e + c f} &
      \frac{b}{a d + b e + c f} &
      \frac{c}{a d + b e + c f}
      \end{array}
      \right)
    \]

    První řádek pravé strany jsou parametry $t_1$, druhý řádek parametry $t_2$.

    \begin{align*}
      A \cdot \left(\begin{array}{c} 1 \\ 0 \\ 0 \end{array}\right)
        = t_1 \left(\begin{array}{c} -\frac{b}{a}\\1\\0 \end{array}\right) + t_2 \left(\begin{array}{c} -\frac{c}{a}\\0\\1 \end{array}\right) \\
      A \cdot \left(\begin{array}{c} 0 \\ 1 \\ 0 \end{array}\right)
        = t_1 \left(\begin{array}{c} -\frac{b}{a}\\1\\0 \end{array}\right) + t_2 \left(\begin{array}{c} -\frac{c}{a}\\0\\1 \end{array}\right) \\
      \underbrace{A \cdot \left(\begin{array}{c} 0 \\ 0 \\ 1 \end{array}\right)
        = t_1 \left(\begin{array}{c} -\frac{b}{a}\\1\\0 \end{array}\right) + t_2 \left(\begin{array}{c} -\frac{c}{a}\\0\\1 \end{array}\right) }
        _{\text{Vypočtěme najednou vynásobením maticí}}
    \end{align*}
    \begin{align*}
        A = A \cdot \left(\begin{array}{c} 1 \\ 0 \\ 0 \end{array}\right)+
          A \cdot \left(\begin{array}{c} 0 \\ 1 \\ 0 \end{array}\right)+
          A \cdot \left(\begin{array}{c} 0 \\ 0 \\ 1 \end{array}\right)
          =
        \underbrace{\begin{pmatrix}
        -\dfrac{b}{a} & -\dfrac{c}{a} \\
        1 & 0 \\
        0 & 1
        \end{pmatrix}}
        _{\text{vektory roviny}}
        \underbrace{\begin{pmatrix}
        -\dfrac{a e}{a d + b e + c f} &
        \dfrac{a d + c f}{a d + b e + c f} &
        -\dfrac{c e}{a d + b e + c f} \\
        -\dfrac{a f}{a d + b e + c f} &
        -\dfrac{b f}{a d + b e + c f} &
        \dfrac{a d + b e}{a d + b e + c f}
        \end{pmatrix}}
        _{\text{z prvních dvou řádků (}t_1,t_2\text{) eliminované matice výše}}
    \end{align*}
    Z matice vpravo vytkněme skalár $\dfrac{1}{ad+be+cf}$, součin matic zadejme do Wolframu a výsledná matice je:
    \begin{align*}
        A = 
        \dfrac{\begin{pmatrix}
        b e + c f & -b d & -c d \\
        -a e & a d + c f & - c e \\
        -a f & - b f & a d + b e
          \end{pmatrix}}{
            ad+be+cf
          }
    \end{align*}

    Zde vidíme smysluplnou podmínku $ad+be+cf \neq 0$,
    která má geometrický význam, že rovina s přímkou ze zadání nesmí být rovnoběžné.

\end{document}
